% Ad Familiares 14.17 --- headnote
% ad-familiares-14-17, approximately 22 December 48 BC, Brundisium

\textit{Cicero to his wife Terentia, written from
Brundisium in the second half of December 48 BC. The
manuscript dateline (\textit{Scr. Brundisi a. 706 (48)
circ. xl i K. Ian.}) is corrupt --- ``xl i'' is
impossible as a Roman date numeral --- and is most
plausibly read as \textit{circ. xi K. Ian.}
(approximately 22 December) or \textit{circ. xvi K. Ian.}
(approximately 17 December); either way, the letter
belongs to the depths of Cicero's Brundisium winter,
the same dark window as the great catastrophe-letters
to Atticus (\textit{Att.} 11.7--11.9).}

\textit{The letter is six lines. Pharsalus has been
lost; Pompey is dead in Egypt; Cicero, having accepted
Caesar's permission to return to Italy, is now stranded
at Brundisium without the political cover of a return
to Rome, estranged from his brother Quintus, broken in
spirit, and unable to bring himself to write at length
even to his wife. The standard Roman epistolary opening
--- \textit{si vales, habetur est, valeo}, ``if you are
well, it is well; I am well'' --- is here all but the
whole letter. He has nothing to say. Lepta and Trebatius,
the bearers, can describe his condition; he can only
ask her to mind her own health and Tullia's. The silence
of the page is the content.}
